\begin{proof}
Razonemos por inducción sobre $n=\dim(C)$. El lema es cierto por hipótesis
si $n=0$, pues entonces $C=k$. Sea $t$ un elemento del ideal maximal
$\mathfrak m$ de $C$ que no pertenece a $\mathfrak m^2$; se sabe entonces que
$C/tC$ es regular y que $\dim(C/tC)=n-1$ ((\textbf 0, 17.1.8) y (\textbf 0, 16.3.4)).
Como
$(D/tD)\otimes_{C/tC}k=D\otimes_C k$,
y $D/tD\in\boldsymbol{C}$, la hipótesis de inducción muestra que se cumple
$\boldsymbol R(D/tD)$; además, $t$ es $C$-regular y, por planitud,
también $D$-regular (\textbf 0$_{\mathrm I}$, 6.3.4). La hipótesis sobre
$\boldsymbol R$ muestra, por tanto, que se cumple $\boldsymbol R(D)$.
\end{proof}
